\section{Практическая работа №3}
\label{sec:pr3}

\subsection{Выполненные задания}

Практическая работа состояла из двух проектов. В проекте \mintinline{text}{glitcher} была реализована обработка изображения: случайная замена байта, сортировка участка байтов и последовательное применение нескольких действий с сохранением промежуточных файлов. Действия были перенесены в \mintinline{haskell}{main}, чтобы точка входа явно читала файл, применяла операции и записывала результаты.

В проекте \mintinline{text}{maze-rws} был реализован лабиринт и поиск пути. Лабиринт представлен как граф комнат, а поиск выполняется обходом в глубину. Основная функция обхода --- \mintinline{haskell}{tryRooms}: она перебирает соседние комнаты, рекурсивно углубляется в непосещённые варианты и возвращает найденный путь.

Основные коммиты практики: \mintinline{text}{d04a59e}, \mintinline{text}{b3058b5}, \mintinline{text}{f94caad}, \mintinline{text}{5f8aa15}, \mintinline{text}{675192e}, \mintinline{text}{aac32e4}, \mintinline{text}{832e62f}, \mintinline{text}{b766d6e}, \mintinline{text}{5d15e30}, \mintinline{text}{88cc5cd}.

\subsection{Использование нейросетей}

В проекте \mintinline{text}{glitcher} был добавлен файл с историей запросов к LLM. Модель использовалась для обсуждения структуры проекта, разбиения чистых и \mintinline{haskell}{IO}-функций и переноса действий в \mintinline{haskell}{main}.

\subsection{Вопросы защиты и их решения}

На защите обсуждалось, где используется монада \mintinline{haskell}{RWS} и почему функция лабиринта имеет такой тип. \mintinline{haskell}{Reader} хранит неизменяемое окружение, то есть карту лабиринта; \mintinline{haskell}{Writer} накапливает лог обхода; \mintinline{haskell}{State} хранит состояние обхода. Такой тип нужен, потому что поиск одновременно читает лабиринт, меняет состояние и записывает сообщения.

Также обсуждалось, где именно находится обход в глубину. В данной реализации он находится в функции \mintinline{haskell}{tryRooms}: функция берёт список возможных комнат и пробует их по одной, углубляясь рекурсивно до тех пор, пока не найдёт целевую комнату или не переберёт все варианты.

\newpage
